% Focal assembly from I/sections/aritmetica_historias.tex, read in full.
% Provenance: technical/PROCEDENCIA_ARITMETICA_HISTORIAS_VI.md.
% hist:seccion and hist:doble-varianza are retained for the common kernel.
% Does not include the Weyl representation or the exceptional prolongation of I.
\section{Arithmetic of histories, composition, and refinement}
\label{vi:hist:seccion-principal}

The persistence of a particle and the evaluation of its character require
a distinction between the transported history and the coordinates read from it.
The APP--TRIT--TPK kernel provides the admissible states, their routes, and
the operations preserving their provenance. This section specifies three
structures used in passing from those data to the atlas: compatibility
between depths, composition of histories with memory, and the
transport of observables in the direction opposite to refinement.
Algebraic coefficients act on states and routes already specified;
the relations proved below preserve that order.

\subsection{Compatible sections and evaluation of their readers}
\label{hist:seccion}

Let \(X_n\) be the set of admissible states at depth \(n\)
produced by the kernel. A state records the APP positions, the
additive and multiplicative sheets, their remainders and quotients, the tritic
regime, orientation, phase, boundary, route, carry, and
memory pertinent to survival. These data satisfy the
relations of the operations that produced them; they are not regarded as a
product of independent coordinates.

For \(m\ge n\), the restriction
\(r_{m,n}:X_m\longrightarrow X_n\) preserves the antecedent at depth
\(n\). The restrictions satisfy
\begin{equation}
 r_{n,n}=\operatorname{id}_{X_n},\qquad
 r_{\ell,n}=r_{m,n}\circ r_{\ell,m}
 \quad(\ell\ge m\ge n).
 \label{vi:hist:restricciones}
\end{equation}
New memory belongs to the prolongation; restriction preserves the memory
already determined in the prefix.

\begin{definition}[Compatible numerical section]
An HMT numerical section is a sequence of states of the same genealogy
\begin{equation}
 x=(x_n)_{n\ge0}\in X_\infty:=\varprojlim_n X_n,
 \qquad r_{m,n}(x_m)=x_n\quad(m\ge n).
 \label{vi:hist:limite-inverso}
\end{equation}
A reader is an operation defined on a specified domain of these
sections. A finite observation belongs to a prefix; a scalar value
is the evaluation of a reader. The section and the presentation of that section
through a reader are therefore objects of different types.
\end{definition}

Inverse-limit notation gathers together the compatibility of the states already
constructed. The existence of a particular prolongation is established
through the admissibility rules of that construction, not merely by
writing \(X_\infty\).

In the Archimedean domain of a reader, each prefix determines a
nonempty closed interval \(I_n(x)=[a_n(x),b_n(x)]\), with
\begin{equation}
 I_{n+1}(x)\subseteq I_n(x),\qquad
 b_n(x)-a_n(x)\longrightarrow0.
 \label{vi:hist:cilindros}
\end{equation}
Regional interval evaluation and the determination of its prefixes
have been assembled in Proposition~\ref{vi:dep:prefijos}.

\begin{proposition}[Value determined by compatible cylinders]
For every section in that domain there is a unique real number
\(\operatorname{ev}(x)\) such that
\begin{equation}
 \{\operatorname{ev}(x)\}=\bigcap_{n\ge0}I_n(x).
 \label{vi:hist:evaluacion}
\end{equation}
\end{proposition}
\begin{proof}
Nestedness makes the sequence \(a_n(x)\) increasing and
\(b_n(x)\) decreasing. Both remain bounded in the initial interval. By
completeness of the line, \(a=\sup_n a_n(x)\) and
\(b=\inf_n b_n(x)\) exist and satisfy \(a\le b\).
The vanishing width forces \(a=b\). This point belongs to every
interval, and any other point of the intersection would have distance from
it less than or equal to every width, and would therefore coincide with it.
\end{proof}

This proposition determines the value of each section in the stated domain.
Coverage of all of \(\mathbb R\) requires a proof that each value
belongs to the image of the reader; the inverse-limit definition does not
by itself provide that proof. Likewise, identifying two
histories from their value requires an injectivity result
for the reader in question. Evaluation preserves the genealogy as
an antecedent and specifies exactly which part of it is expressed.

\subsection{Sectorial route algebra and memory multipliers}
\label{vi:hist:algebra}

Fix a sector of states and the category \(\mathcal C\) of its
admissible histories. We write \(vu=v\circ u\): first
\(u\) is traversed, then \(v\). The elementary steps are TPK
selection, transport, and update operations. The category preserves
their endpoints and the conditions allowing two routes to be concatenated.

The local tritic fibre can be represented by
\[
 \mathbb R[J_x]/(J_x^2+\tau(x)1_x),
 \qquad \tau(x)\in\{+1,0,-1\}.
\]
The representation preserves the regime and its orientation. To handle
richer coefficients, we shall allow an associative unital real
algebra \(A_x\) at each object \(x\). A transition between regimes
must specify its map on these algebras; it does not amount to identifying them.

The following construction applies to sectors possessing unital
transports \(\rho_u:A_x\to A_y\), for \(u:x\to y\), with
\begin{equation}
 \rho_{\operatorname{id}_x}=\operatorname{id}_{A_x},\qquad
 \rho_v\circ\rho_u=\rho_{vu}.
 \label{vi:hist:accion}
\end{equation}
For each composable pair \(u:x\to y\), \(v:y\to z\), let
\(\Omega(v,u)\in Z(A_z)^\times\) be an invertible central multiplier.
Its rules are
\begin{equation}
 \Omega(\operatorname{id}_y,u)
 =\Omega(u,\operatorname{id}_x)=1_y
 \label{vi:hist:normalizacion}
\end{equation}
and, for \(w:z\to t\),
\begin{equation}
 \rho_w\!\left(\Omega(v,u)\right)\Omega(w,vu)
 =\Omega(w,v)\Omega(wv,u).
 \label{vi:hist:cociclo}
\end{equation}
These conditions describe the starting structure of the sector: every
specific application must provide its transports and multipliers.
A realisation through bimodules preserves its own compositions and
is not automatically replaced by the strict action
\eqref{vi:hist:accion}.

Let \(t(u)\) be the terminal endpoint of the arrow \(u\). On sums
of finite support
\[
 \mathcal H=\bigoplus_{u\in\operatorname{Mor}(\mathcal C)}A_{t(u)}T_u
\]
we define, by bilinearity,
\begin{equation}
 (aT_v)(bT_u)=
 a\rho_v(b)\Omega(v,u)T_{vu}
 \label{vi:hist:producto}
\end{equation}
when \(u,v\) are composable, and zero otherwise. Addition gathers
histories with their coefficients; multiplication concatenates the histories and
preserves the memory multiplier. Finite support ensures that each
product is defined without requiring additional convergence.

\subsection{Associativity, local identities, and calendar memory}
\label{vi:hist:asociatividad}

\begin{proposition}[Associative composition of the sector]
Under \eqref{vi:hist:accion}--\eqref{vi:hist:cociclo}, the product
\eqref{vi:hist:producto} is associative. If there are finitely many objects, its identity
is \(\sum_x1_xT_{\operatorname{id}_x}\). For an arbitrary set of objects,
the same expressions over finite sets of endpoints
give local identities for every finite collection of histories.
\end{proposition}
\begin{proof}
Let \(u:x\to y\), \(v:y\to z\), \(w:z\to t\), with coefficients
\(a\in A_y\), \(b\in A_z\), \(c\in A_t\). The two associations
produce, as coefficients of \(T_{wvu}\),
\begin{align*}
 ((cT_w)(bT_v))(aT_u):\quad&
 c\rho_w(b)\Omega(w,v)\rho_{wv}(a)\Omega(wv,u),\\
 (cT_w)((bT_v)(aT_u)):\quad&
 c\rho_w(b)\rho_w\rho_v(a)
 \rho_w(\Omega(v,u))\Omega(w,vu).
\end{align*}
The strict action identifies \(\rho_w\rho_v(a)\) with
\(\rho_{wv}(a)\). Centrality allows
\(\Omega(w,v)\) to be moved to the right of this factor, and the cocycle identity
equates the remaining factors. The general coefficients
\(a,b,c\), which need not commute, are not permuted. If the three arrows do not form a
composable chain, both associations are zero by their endpoints.
Bilinearity extends the equality to finite sums. Finally,
\eqref{vi:hist:normalizacion} and the unitality of each \(\rho_u\)
verify the identity. For a finite collection it suffices to add the identities
of all its initial and terminal endpoints, proving the last
assertion.
\end{proof}

The nonadic calendar provides an effective realisation of these
conditions. Take the one-object category with arrows
\(C_9=\mathbb Z/9\mathbb Z\), coefficients
\(A=\mathbb R[z,z^{-1}]\), and identity action. For representatives
\(0\le a,b<9\), define
\begin{equation}
 c_0(a,b)=\left\lfloor\frac{a+b}{9}\right\rfloor,
 \qquad\Omega(b,a)=z^{c_0(a,b)}.
 \label{vi:hist:carry-calendario}
\end{equation}
The formal variable \(z\) preserves the number of circuits. Normalisation
follows from \(c_0(0,a)=c_0(a,0)=0\). If \([a+b]_9\) denotes the
reduced representative, Euclidean division gives
\[
 c_0(a,b)+c_0([a+b]_9,c)
 =\left\lfloor\frac{a+b+c}{9}\right\rfloor
 =c_0(b,c)+c_0(a,[b+c]_9).
\]
This equality proves the cocycle \eqref{vi:hist:cociclo} for the
multipliers, and hence the associativity of the calendar with memory.

If \(T=T_{[1]_9}\), the product rules give
\(T^j=T_{[j]_9}\) for \(0\le j<9\) and \(T^9=zT_{[0]_9}\).
Subsequently imposing \(z^{12}=1\) yields \(T^{108}=1\).
The expressions \(z^kT_{[j]_9}\), \(0\le k<12\), \(0\le j<9\),
form the \(108\) positions of the resulting calendar. This recovers
the cyclic group \(C_{108}\) and its projection onto \(C_9\); retaining
\(z\) without that relation distinguishes successive circuits. This is a
realisation of calendar memory: the remaining coordinates of
the TPK histories remain in their own domain.

\subsection{Double variance of refinement and preservation of unity}
\label{hist:doble-varianza}

Restriction of states induces a prolongation of observables in the opposite
direction. For the function algebras with pointwise multiplication
\(D_n=\operatorname{Fun}(X_n,\mathbb R)\), precomposition defines
\begin{equation}
 r_{m,n}^*:D_n\longrightarrow D_m,
 \qquad r_{m,n}^*f=f\circ r_{m,n}.
 \label{vi:hist:precomposicion}
\end{equation}
The contravariant order is preserved:
\(r_{\ell,n}^*=r_{\ell,m}^*\circ r_{m,n}^*\). These maps
are unital homomorphisms. If \(r_{m,n}\) is surjective, then
\(r_{m,n}^*\) is also injective: a difference between two functions can be
evaluated at a preimage of a point where they differ.

\begin{lemma}[Indicators, unity, and compatible evaluation]
For \(x\in X_n\), let \(e_x\) be its indicator function. Then
\begin{equation}
 r_{m,n}^*e_x=\mathbf1_{r_{m,n}^{-1}(x)},\qquad
 r_{m,n}^*1_n=1_m.
 \label{vi:hist:unidad}
\end{equation}
When the fibre \(r_{m,n}^{-1}(x)\) is finite, the first equality is
the finite sum
\[
 r_{m,n}^*e_x=\sum_{y:\,r_{m,n}(y)=x}e_y.
\]
For every section \(x=(x_n)\in X_\infty\), the rule
\([f]\mapsto f(x_n)\), with \(f\in D_n\), defines a unital
homomorphism
\begin{equation}
 \operatorname{Ev}_x:D_{\mathrm{cyl}}\longrightarrow\mathbb R,
 \qquad
 D_{\mathrm{cyl}}=\varinjlim_n(D_n,r_{n+1,n}^*).
 \label{vi:hist:evaluacion-observables}
\end{equation}
\end{lemma}
\begin{proof}
For \(y\in X_m\), we have
\((r_{m,n}^*e_x)(y)=1\) exactly when \(r_{m,n}(y)=x\).
This is the first equality of \eqref{vi:hist:unidad}. If the fibre is
finite, its point indicators are orthogonal and sum to its indicator.
The second equality follows from \(1_n(r_{m,n}(y))=1\), without any
assumption on the cardinality of the fibre.

For every \(f\in D_n\), the naturality of evaluation is
\[
 (r_{m,n}^*f)(y)=f(r_{m,n}(y)).
\]
In a compatible section, \(r_{m,n}(x_m)=x_n\), hence
\((r_{m,n}^*f)(x_m)=f(x_n)\). Thus two representatives identified
at a common depth have the same evaluation.
Point evaluation preserves addition, multiplication, and unity, completing
the proof of \eqref{vi:hist:evaluacion-observables}.
\end{proof}

For infinite fibres the formulation through
\(\mathbf1_{r_{m,n}^{-1}(x)}\) remains valid; there is no need to interpret an
infinite sum as a finite algebraic sum. The direct limit gathers the
finite-depth observables with the identifications induced by
the restrictions. Prolonging route operators additionally requires
the compatibility squares of their transports with those restrictions.

The preserved unity is the constant function one on the domain of
possibilities. An individual section selects a compatible history
in that domain. Refinement may increase depth and distinguish
more prolongations while preserving unity and the evaluation of each
antecedent. This unital preservation provides the precise algebraic content
used here; a probability measure, when it enters
another reader, additionally requires its normalisation and its own rules.
